% =====================================================================
% LAMPIRAN I: KUESIONER EVALUASI USABILITAS (SUS) SMARTCITYAPPS
% =====================================================================
\chapter{Kuesioner Evaluasi Usabilitas SmartCityApps}
\label{app:kuesioner-sus}

Lampiran ini memuat formulir lengkap evaluasi pengalaman pengguna yang dibagikan kepada responden Warga Pelapor. Formulir terdiri atas lima bagian: profil responden, pakta integritas, skenario tugas yang dijalankan responden, kesepuluh pernyataan \textit{System Usability Scale} (SUS), dan pertanyaan umpan balik kualitatif. Kesepuluh pernyataan SUS menggunakan terjemahan Bahasa Indonesia yang telah divalidasi oleh \citet{sharfinaIndonesianAdaptationSystem2016}. Metodologi penghitungan dan interpretasi skornya dijelaskan pada Bagian~\ref{sec:metode-usabilitas}.

\section*{Bagian A: Profil Responden}

\begin{table}[htbp]
\centering
\small
\caption{Isian profil responden.}
\label{tab:form-profil}
\begin{tabular}{p{4.5cm} p{8cm}}
\toprule
\textbf{Atribut} & \textbf{Isian} \\
\midrule
Nama / Inisial & \rule{7.5cm}{0.4pt} \\
\addlinespace
Usia & \rule{7.5cm}{0.4pt} \\
\addlinespace
Jenis kelamin & $\square$ Laki-laki \quad $\square$ Perempuan \\
\addlinespace
Pekerjaan / Status & \rule{7.5cm}{0.4pt} \\
\addlinespace
Frekuensi penggunaan \textit{smartphone} & $\square$ Jarang \quad $\square$ Kadang-kadang \quad $\square$ Sering \\
\addlinespace
Pernah menggunakan aplikasi pelaporan publik (mis. LAPOR!/JAKI) & $\square$ Pernah \quad $\square$ Belum pernah \\
\bottomrule
\end{tabular}
\end{table}

\section*{Bagian B: Pakta Integritas}

Dengan ini saya menyatakan bahwa saya bersedia menjadi responden pada evaluasi usabilitas aplikasi SmartCityApps. Saya mengisi kuesioner ini secara sukarela, jujur, dan berdasarkan pengalaman saya sendiri saat mencoba aplikasi. Saya memahami bahwa jawaban saya digunakan semata-mata untuk keperluan penelitian tugas akhir dan kerahasiaan identitas saya dijaga.

\vspace{0.8cm}
\noindent Tanggal: \rule{3.5cm}{0.4pt} \hfill Tanda tangan: \rule{4cm}{0.4pt}

\section*{Bagian C: Skenario Tugas}

Sebelum mengisi penilaian, responden diminta mencoba menyelesaikan skenario tugas berikut menggunakan aplikasi SmartCityApps pada perangkat Android. Pendamping pengujian mencatat keberhasilan tiap tugas.

\begin{table}[htbp]
\centering
\small
\caption{Skenario tugas yang dijalankan responden.}
\label{tab:form-skenario}
\begin{tabular}{p{1.2cm} p{8.5cm} p{2.4cm}}
\toprule
\textbf{Kode} & \textbf{Tugas} & \textbf{Berhasil?} \\
\midrule
T-01 & Mendaftar akun baru atau masuk ke akun yang sudah ada. & $\square$ Ya \; $\square$ Tidak \\
T-02 & Membuka layar buat laporan dan mengambil foto kejadian dengan kamera. & $\square$ Ya \; $\square$ Tidak \\
T-03 & Memilih foto kejadian dari galeri perangkat. & $\square$ Ya \; $\square$ Tidak \\
T-04 & Mengisi narasi atau deskripsi laporan. & $\square$ Ya \; $\square$ Tidak \\
T-05 & Menetapkan lokasi laporan melalui GPS atau memilih titik di peta. & $\square$ Ya \; $\square$ Tidak \\
T-06 & Menjalankan klasifikasi AI dan meninjau hasil prediksi instansi, \textit{confidence}, serta instansi terdekat. & $\square$ Ya \; $\square$ Tidak \\
T-07 & Mengoreksi kategori instansi secara manual pada layar tinjauan AI. & $\square$ Ya \; $\square$ Tidak \\
T-08 & Mengirim laporan lalu memeriksa kemunculannya beserta status pada riwayat. & $\square$ Ya \; $\square$ Tidak \\
\bottomrule
\end{tabular}
\end{table}

\section*{Bagian D: Pernyataan SUS}

Berikan penilaian terhadap setiap pernyataan dengan memberi tanda pada salah satu kolom skala, dengan ketentuan: 1 = Sangat Tidak Setuju (STS), 2 = Tidak Setuju (TS), 3 = Netral (N), 4 = Setuju (S), dan 5 = Sangat Setuju (SS).

\begin{table}[htbp]
\centering
\small
\caption{Kesepuluh pernyataan SUS dengan skala Likert lima poin.}
\label{tab:form-sus}
\begin{tabular}{c p{7.2cm} c c c c c}
\toprule
\textbf{No.} & \textbf{Pernyataan} & \textbf{1} & \textbf{2} & \textbf{3} & \textbf{4} & \textbf{5} \\
 & & \textbf{STS} & \textbf{TS} & \textbf{N} & \textbf{S} & \textbf{SS} \\
\midrule
P1 & Saya berpikir akan menggunakan aplikasi ini lagi. & & & & & \\
\addlinespace
P2 & Saya merasa aplikasi ini rumit untuk digunakan. & & & & & \\
\addlinespace
P3 & Saya merasa aplikasi ini mudah digunakan. & & & & & \\
\addlinespace
P4 & Saya membutuhkan bantuan dari orang lain atau teknisi dalam menggunakan aplikasi ini. & & & & & \\
\addlinespace
P5 & Saya merasa fitur-fitur aplikasi ini berjalan dengan semestinya. & & & & & \\
\addlinespace
P6 & Saya merasa ada banyak hal yang tidak konsisten (tidak serasi) pada aplikasi ini. & & & & & \\
\addlinespace
P7 & Saya merasa orang lain akan memahami cara menggunakan aplikasi ini dengan cepat. & & & & & \\
\addlinespace
P8 & Saya merasa aplikasi ini membingungkan. & & & & & \\
\addlinespace
P9 & Saya merasa tidak ada hambatan dalam menggunakan aplikasi ini. & & & & & \\
\addlinespace
P10 & Saya perlu membiasakan diri terlebih dahulu sebelum menggunakan aplikasi ini. & & & & & \\
\bottomrule
\end{tabular}
\end{table}

\section*{Bagian E: Umpan Balik Kualitatif}

\noindent 1. Menurut Anda, apa kelebihan utama aplikasi SmartCityApps?

\vspace{0.2cm}
\noindent \rule{\linewidth}{0.4pt}
\vspace{0.2cm}
\noindent \rule{\linewidth}{0.4pt}

\vspace{0.4cm}
\noindent 2. Bagian atau fitur apa yang menurut Anda masih membingungkan atau perlu diperbaiki?

\vspace{0.2cm}
\noindent \rule{\linewidth}{0.4pt}
\vspace{0.2cm}
\noindent \rule{\linewidth}{0.4pt}

\vspace{0.4cm}
\noindent 3. Apakah Anda bersedia menggunakan aplikasi ini untuk melaporkan masalah lingkungan di sekitar Anda? Mohon jelaskan alasannya.

\vspace{0.2cm}
\noindent \rule{\linewidth}{0.4pt}
\vspace{0.2cm}
\noindent \rule{\linewidth}{0.4pt}
